\documentclass[11pt,a4paper]{article}
\usepackage[margin=0.85in,top=1in,bottom=1in]{geometry}
\usepackage{amsmath,amssymb,amsfonts}
\usepackage{graphicx}
\usepackage{float}
\usepackage{enumitem}
\usepackage{microtype}
\usepackage{caption}
\usepackage[hidelinks]{hyperref}
\usepackage[most,breakable]{tcolorbox}
\tcbuselibrary{breakable,skins}
\usepackage{fontspec}
\setmainfont{DejaVu Serif}
\setsansfont{DejaVu Sans}
\setmonofont{DejaVu Sans Mono}
\captionsetup{font=small,labelfont=bf,skip=4pt}
\setlist{leftmargin=1.5em,topsep=2pt}

%% Breakable problem card — prevents content cutoff across pages
\newtcolorbox{solcard}[3]{
  breakable, enhanced,
  colback=white, colframe=gray!50, arc=2pt,
  title={\textbf{#1}\hfill\textsf{\small #3\,pts}},
  fonttitle=\normalsize\sffamily\bfseries,
  before upper={\small\textit{#2}\par\smallskip\hrule height 0.4pt\smallskip},
  top=4pt, bottom=6pt, left=7pt, right=7pt,
  parbox=false,
}

\title{\textbf{Physics of rubber --- Solution}}
\author{\normalsize Y26 (2026) — English translation}
\date{}
\begin{document}\maketitle\vspace{-0.5em}
\begin{solcard}{A1}{Find the number of microscopic states $\Gamma$, with the help of which the macroscopic state is realized, in which the difference in the number of links oriented along and against the axis $x$ is equal to $m = n_1 - n_2$. Also obtain an approximate answer for $m \ll n$. Express your answers using $m$ and $n$.}{0.40}

For a given $m$, the number of links oriented in the positive direction is $n_1 = \dfrac{n+m}{2}$, and in the negative direction $n_2 = \dfrac{n-m}{2}$. Then the required number of ways

\par\smallskip\noindent\textbf{Answer:} \textbackslash{}textbf{Answer:} $$ \Gamma = \frac{n!}{\left(\frac{n+m}{2}\right)! \left(\frac{n-m}{2}\right)!} \approx 2^n \sqrt{\frac{2}{\pi n}} \exp\left( - \frac{m^2}{2n}\right) $$
\end{solcard}

\begin{solcard}{A2}{Show that at $ h_x \ll n a$ the entropy of the molecule has the form $$ S = S_0 - \gamma h_x^2. $$Express coefficient $\gamma$ via/through $n$, $a$ and $k$.}{0.80}

The displacement of the end of the molecule is $h_x = (n_1 - n_2) a = ma$. Then, using the formula from the previous paragraph, we get $$ \Gamma = 2^n \sqrt{\frac{2}{\pi n}} \exp\left( - \frac{h_x^2}{2n a^2}\right). $$Then from the definition энтропии we get $$ S = k \ln \Gamma = k \ln\left( 2^n \sqrt{\frac{2}{\pi n}} \right) - \frac{k h_x^2}{2 n a^2}. $$

\par\smallskip\noindent\textbf{Answer:} \textbackslash{}textbf{Answer:} $$ \gamma = \frac{k}{2 na^2} $$
\end{solcard}

\begin{solcard}{A3}{The probability that the end of a molecule lies in the interval $[h_x, h_x + d h_x]$ of width $dh_x \gg a$ is $f(h_x) dh_x$. Show that for $|h_x| \ll na$ the probability density of $f(h_x)$ has the form $$ f(h_x) = A \exp(-\alpha h_x^2) $$ and выtimesandте постоянные $A$ and $\alpha$ via/through $n$ and $a$. Считайте, что направления звеньев цепи независимы друг от друга and equalsвероятны. For/Atмечание: for/at повороте одного звена $h_x$ меняется на $2a$.}{0.60}

Note: When one link is rotated, $h_x$ changes to $2a$.


The total number of states in the circuit is $2^n$. Since the directions of the chain links are equally probable, the probability of a given state of the molecule is proportional to the number of microscopic states in which it can be realized. Then the probability that the displacement of the end of the molecule takes on a given value $h_x$ is equal to $$ p(h_x) = \frac{1}{2^n} \Gamma(h_x) = \sqrt{\frac{2}{\pi n}}\exp\left( - \frac{h_x^2}{2n a^2}\right). $$For/At повороте одного звена value $m$ меняется на $2$, therefore смещение $h_x$ меняется на $2 a$. Значит в некотором интервале ширины $dh_x \gg a$ находится $dN = dh_x/2a$ возможных состояний. For/At достаточно малом $dh_x$ вероятности этих состояний практически одинаковы, and therefore вероятность попадания в данный интервал equals $$ dN \cdot p(h_x) = \frac{dh_x}{2 a} \sqrt{\frac{2}{\pi n}}\exp\left( - \frac{h_x^2}{2n a^2}\right). $$Hence density вероятности $$ f(h_x) = \frac{1}{a \sqrt{2 \pi n }}\exp\left( - \frac{h_x^2}{2n a^2}\right). $$Taking into account expression for гауссова интеграла $$ \int\limits_{-\infty}^{+\infty} dx\, e^{- \beta x^2} = \sqrt{\frac{\pi}{\beta}} $$we get $$ \int\limits_{-\infty}^{+\infty} dh_x\, f(h_x) =  \frac{1}{a \sqrt{2 \pi n }} \sqrt{2\pi n a^2} = 1, $$that is, the distribution is correctly normalized.

\par\smallskip\noindent\textbf{Answer:} \textbackslash{}textbf{Answer:} $$ A = \frac{1}{a \sqrt{2 \pi n}}, \quad \alpha = \frac{1}{2 n a^2} $$
\end{solcard}

\begin{solcard}{B1}{What are the values ​​of $\langle h^2\rangle $ and $\langle h_x^2\rangle $ for a three-dimensional chain? Express your answers using $n$ and $l$. Note: if you were unable to complete this step, you can continue to count $\langle h_x^2\rangle = nl^2$ throughout the entire task.}{0.50}

Note: if you were unable to complete this step, you can continue to count $\langle h_x^2\rangle = nl^2$ throughout the entire task.


The displacement vector of the end of the vector relative to the beginning is equal to the sum of the vectors $\vec{l}_i$, $$ \vec{h} = \sum_i \vec{l}_i, $$therefore его квадрат $$ \vec{h}{}^2 = \left( \sum_i \vec{l}_i \right)^2 = \sum_i (\vec{l}_i)^2 +2 \sum_{i < j} \vec{l}_i {\vec{l}_j}. $$Here в первой сумме суммирование производится по всем звеньям молекулы, а во второй – по всем vaporам звеньев. Каждое term первой суммы equals длине одного звена $l^2$, therefore вся первая сумма equals $n l^2$. Усредним вторую сумму по сличным состояниям молекулы. В силу независимости направлений сличных звеньев каждое from произведений for/at усреднении gives 0, $\langle \vec{l}_i \vec{l}_j \rangle = 0$, therefore and вся сумма for/at усреднении обращается в 0. Then for среднего квадрата we get $$ \langle h^2 \rangle = n l^2. $$В силу equalsвероятности всех направлений $$ \langle h_x^2 \rangle = \langle h_y^2 \rangle =\langle h_z^2 \rangle. $$Since $$ \langle h^2 \rangle = \langle h_x^2 \rangle + \langle h_y^2 \rangle+ \langle h_z^2 \rangle, $$средний квадрат одной from проекций equals $$ \langle h_x^2 \rangle = \frac{1}{3} \langle h^2 \rangle = \frac{1}{3} n l^2. $$

\par\smallskip\noindent\textbf{Answer:} \textbackslash{}textbf{Answer:} $$ \langle h^2 \rangle = n l^2 ,\quad \langle h_x^2 \rangle = \frac{1}{3} nl^2 $$
\end{solcard}

\begin{solcard}{B2}{What is the average distance between the ends of the chain $\langle h \rangle$? Express your answer in terms of $n$ and $l$.}{1.00}

Differentiating the value of the Gaussian integral $$ \int\limits_{-\infty}^{+\infty}dx\, e^{-\beta x^2} = \sqrt{\frac{\pi}{\beta}} $$по parameterу $\beta$ we get $$ \int\limits_{-\infty}^{+\infty}dx\, x^2 e^{-\beta x^2} = \frac{1}{2}\sqrt{\frac{\pi}{\beta^3}}, \quad \int\limits_{-\infty}^{+\infty}dx\, x^4 e^{-\beta x^2} =\frac{3}{4} \sqrt{\frac{\pi}{\beta^5}}. $$From the condition нормировки плотности вероятности $$ \int d^3 h f(\vec{h}) = 1 = A \int d^3 h e^{- \beta h^2} = 4 \pi A \int\limits_0^{\infty}h^2 dh\, e^{- \beta h^2} =  \pi A \frac{\sqrt{\pi}}{\beta^{3/2}} $$we find coefficient $$ A = \left(\frac{\beta}{\pi}\right)^{3/2}. $$Note that for/at вычислении интеграла появился дополнительный множитель $1/2$, since нижний limit equals $0$, not $- \infty$. Then for среднего квадрата we get $$ \langle h^2 \rangle = \int d^3 h \, h^2 f(\vec{h}) = 4\pi A \int\limits_0^{\infty}h^4 dh\, e^{-\beta h^2} = \pi A \frac{3}{2} \frac{\sqrt{\pi}}{\beta^{5/2}} = \frac{3}{2 \beta}. $$Using expression from предыдущего partа, we find $$ \frac{3}{2 \beta} = nl^2, $$hence $$ \beta = \frac{3}{2 n l^2}. $$ Let us calculate средний magnitude смещения $$ \langle h \rangle = \int d^3 h\, h f(\vec{h}) = 4\pi A \int\limits_0^{\infty} h^3 dh\, e^{-\beta h^2}. $$Перейдем к переменной $\xi = h^2$: $$ \langle h \rangle = 2\pi A \int\limits_0^{\infty}\xi d \xi\, e^{- \beta \xi} = 2 \pi A\left(\left. - \frac{1}{\beta} \xi e^{-\beta \xi}\right|_0^{\infty} + \frac{1}{\beta}\int\limits_0^{\infty} d \xi\, e^{- \beta \xi}\right) = \frac{2\pi A}{\beta^2} = \frac{2}{\sqrt{\pi \beta}}. $$Substituting value $\beta$, we finally find


The same result can be obtained using an analogy with the Maxwell distribution. Indeed, the distribution of the displacement vector of the second vector of the end of the molecule relative to the first differs from the distribution of the velocity vector only in the value of the parameter $\beta$. For the Maxwell distribution, the root mean square value of the velocity $$ \langle v^2 \rangle  = \frac{3 k T}{m}, $$а also value среднего модуля скорости $$ \langle v \rangle = \sqrt{\frac{8 k T}{\pi m}}. $$Hence we get безсмерное ratio $$ \frac{\langle v \rangle ^2}{\langle v^2 \rangle} = \frac{8}{3\pi}. $$Оно не содержит parameterов расlimitения, therefore такое же соratio должно выполняться and for  $\vec{h}$: $$ \frac{\langle h \rangle ^2}{\langle h^2 \rangle} = \frac{8}{3\pi}, $$ is known, from which we immediately obtain the above answer.

\par\smallskip\noindent\textbf{Answer:} \textbackslash{}textbf{Answer:} $$ \langle h \rangle = \sqrt{\frac{8 n}{3\pi}} l $$
\end{solcard}

\begin{solcard}{B3}{Find the average displacement of the second end of the molecule $h$. Express your answer in terms of $F$, $T$, $n$, $l$ and $k$. Obtain an exact expression and an approximate one for small $F$.}{1.50}

From symmetry, the average displacement vector of the end of the molecule will be directed along the external force, so you only need to find the projection of this vector onto the direction of the force. Let's consider some link, let it form an angle $\theta$ with the direction of the force. Then its potential energy $U = - F l \cos \theta$, which means the probability that it is directed at some solid angle $d \Omega = 2\pi \sin \theta d \theta$ depends only on $\theta$ and has the form $$ B \exp\left(-\frac{U}{kT}\right)d\Omega = B \exp\left(\frac{F l}{kT} \cos \theta \right)d\Omega, $$where coefficient $B$ можно найти from нормировки. For удобства вычислений введем безсмерную переменную $\xi = \dfrac{Fl}{kT}$. Then $$ B \int \limits_0^{\pi} e^{\xi \cos \theta} 2\pi \sin \theta d \theta = 1. $$Перейдем к переменной интегрирования $u = \cos \theta$, которая меняется от $-1$ до $1$, we get $$ 2\pi B \int\limits_{-1}^{1} e^{\xi u} du = \frac{2 \pi B}{\xi} (e^{\xi} - e^{i\xi}) = \frac{4 \pi B}{\xi} \sinh \xi = 1, $$hence $$ B  = \frac{1}{4\pi} \frac{\xi}{\sinh \xi}. $$Then средняя проекция звена на направление силы $$ \langle l \cos \theta \rangle = B \int \limits_0^{\pi} l \cos \theta e^{\xi \cos \theta} 2\pi \sin \theta d \theta = 2\pi B l\int\limits_{-1}^1 udu_\, e^{\xi u}. $⟦1 8⟧$ \int\limits_{-1}^1 udu_\, e^{\xi u} = \frac{1}{\xi}\left. ue^{\xi u}\right|_{-1}^1 -\frac{1}{\xi} \int\limits_{-1}^1du\, e^{\xi u} = \frac{1}{\xi} (e^{\xi u} + e^{-\xi u}) - \frac{1}{\xi^2} (e^{\xi u} - e^{- \xi u}). $$$$ \langle l \cos \theta \rangle = \frac{l}{2} \frac{\xi}{\sinh \xi} 2\left(\frac{\cosh \xi}{\xi} - \frac{\sinh \xi}{\xi^2} \right) = l \left( \coth \xi - \frac{1}{\xi}\right). $$Полное смещение конца молекул складывается from проекций отдельных звеньев: $$ h = n \langle l \cos \theta \rangle = n l \left(\coth \frac{Fl}{kT} - \frac{kT}{Fl}\right).$$В limitе малой силы we get сложение $$ \coth \xi - \frac{1}{\xi} = \frac{e^{\xi} + e^{-\xi}}{e^{\xi}-e^{-\xi}} - \frac{1}{\xi} \approx \frac{2 + \xi^2}{2 \xi + \xi^3/3}- \frac{1}{\xi} = \frac{1}{\xi}\left( \frac{1+ \xi^2/2}{1 + \xi^2/6}-1\right) = \frac{1}{\xi} \frac{\xi^2}{3} = \frac{\xi}{3}. $$Then for смещения we get $$ h \approx \frac{nl}{3} \frac{Fl}{kT} = \frac{nl^2}{3 kT}F. $$Note also, что в limitе большой силы, that is $\xi \to \infty$ we get $\coth \xi \to 1$ and $$ h \approx nl \left(1 - \frac{kT}{Fl} \right) \approx nl, $$that is молекула практически полностью выпрямляется. Характерное value силы, for/at котором происходит переход from области малых сил в область больших сил, отвечает $\xi \sim 1$, that is $F \sim kT/l$.

\par\smallskip\noindent\textbf{Answer:} \textbackslash{}textbf{Answer:} $$ h =  n l \left(\coth \frac{Fl}{kT} - \frac{kT}{Fl}\right) \approx \frac{n l^2}{3k T} F$$
\end{solcard}

\begin{solcard}{B4}{Plot a graph of the length of the molecule $h$ versus the applied force. Provide typical values.}{0.50}
\begin{figure}[H]
  \centering
  \includegraphics[width=0.72\textwidth]{T2_S_EN_assets/fig1.jpg}
\end{figure}
\end{solcard}

\begin{solcard}{С1}{What is the change in entropy $\Delta S$ of the entire mesh as a result of stretching? Express your answer in terms of $k,~\lambda_x,~\lambda_y,~\lambda_z$ and $N$. Note: if you were unable to make B1, then use $\langle h_x^2\rangle = nl^2$.}{0.70}

Note: if you were unable to make B1, then use $\langle h_x^2\rangle = nl^2$.


The entropy of one chain in a stretched state is equal to \[ S_1(\vec{h'}) = S_0 - \dfrac{3k}{2nl^2}(\lambda_x^2h_x^2 + \lambda_y^2h_y^2 + \lambda_z^2h_z^2).\]The change in the entropy of one chain as a result of stretching \[\Delta S_1 = S_1(\vec{h'}) - S_1(\vec{h}) = - \dfrac{3k}{2nl^2}((\lambda_x^2-1)h_x^2 + (\lambda_y^2-1)h_y^2 + (\lambda_z^2-1)h_z^2).\]The change in the entropy of the polymer network is expressed through the average value of the change in the entropy of one chain: \[ \Delta S = N\langle\Delta S_1\rangle = - \dfrac{3kN}{2nl^2}((\lambda_x^2-1)\langle h_x^2 \rangle+ (\lambda_y^2-1)\langle h_y^2 \rangle + (\lambda_z^2-1)\langle h_z^2\rangle).\] Using the result of point B, we find

\par\smallskip\noindent\textbf{Answer:} \textbackslash{}textbf{Answer:} \[\Delta S = -\dfrac{kN}{2}(\lambda_x^2 + \lambda_y^2 +\lambda_z^2 -3)\]
\end{solcard}

\begin{solcard}{С2}{Let us consider tension in the case when forces are applied along one of the axes and stretch the sample along this axis by a factor of $\lambda$, and there are no stresses along the other axes. What is the change in entropy of the sample? Express your answer in terms of $k,~\lambda$ and $N$.}{0.60}

If forces are applied only along the $x$ axis, then $\lambda_x = \lambda$, and for other axes $\lambda_y = \lambda_z$ is true. Since the volume of the rod does not change, then $\lambda_x\lambda_y\lambda_z = 1$. Then $\lambda_y = \lambda_z = 1/\sqrt{\lambda}$. Substituting the dilations into the formula from C2, we find

\par\smallskip\noindent\textbf{Answer:} \textbackslash{}textbf{Answer:} \[\Delta S = -\dfrac{kN}{2}\left(\lambda^2 + \dfrac{2}{\lambda}-3\right)\]
\end{solcard}

\begin{solcard}{C3}{Using the first law of thermodynamics for an isothermal process and the result of part C2, show that the equation of state of rubber (the dependence of the elastic force on tension and temperature) is: \[ F(\lambda,T) = \beta(T) \left(\lambda-\dfrac{1}{\lambda^2}\right).\] Find $\beta$. Express your answer in terms of $k,~N,~T$ and the length of the sample in the unstretched state $l_0$.}{0.60}

The first beginning in general form is written as $$ dU = \delta Q + \delta A. $$Here подведенное тепло можно выtimesandть via/through изменение энтропии: $\delta Q = T dS$. Под работой $\delta A$ for/at таком выборе signов подсумевается work, совершенная над обсцом. Она в общем сrayае складывается from работы сил давления and работы внешних  сил, растягивающих резину $$ \delta A = Fdl - pdV. $$ Then we finally get ⟦0⟧where $P$ – внешнее pressure, $dV$ – изменение volumeа стержня. Since volume стержня не изменяется, внешнее pressure не совершает работу. Since внутренняя energy зависит только от температуры, то в изотермическом процессе $dU = 0$. Then \[ F = -T\left(\dfrac{\partial S}{\partial l}\right)_T = -\dfrac{T}{l_0}\left(\dfrac{\partial \Delta S}{\partial \lambda}\right)_T =\dfrac{kNT}{l_0}\left(\lambda - \dfrac{1}{\lambda^2}\right).\]

\par\smallskip\noindent\textbf{Answer:} \textbackslash{}textbf{Answer:} \[\beta = \dfrac{kNT}{l_0}\]
\end{solcard}

\begin{solcard}{C4}{Show that Hooke's law is valid for small deformations. Find the Young's modulus $E$ of the rubber. Express your answer in terms of $k,~T$ and the concentration $\nu$ of polymer chains in the sample. Calculate the numerical value of Young's modulus if the rubber density is $\rho = 1200~\text{kg}/\text{m}^3$, molar mass $\mu = 10^5~\text{g}/\text{mol}$, temperature $T = 300~\text{K}$.}{0.80}

Let us substitute into the expression for force $\lambda = 1 + \varepsilon$ at $\varepsilon \ll 1$: \[F=\dfrac{kNT}{l_0}\left(1+\varepsilon - \dfrac{1}{(1+\varepsilon)^2}\right) \approx \dfrac{3kNT\varepsilon}{l_0}.\] Thus, for small deformations, the force is directly proportional to the relative elongation, i.e. Hooke's law is true. Young's modulus is equal to ($S_0$ – initial cross-sectional area) \[ E = \dfrac{F}{S_0\varepsilon} = \dfrac{3kNT}{l_0 S_0} = \dfrac{3kNT}{V} = 3k\nu T.\]The concentration of polymer chains can be determined by knowing the rubber density and molar mass: $$ \nu = \frac{\rho}{\mu} N_A, $$since $\rho/\mu$ – the number of moles per unit volume. Taking into account the numerical values ​​\[\nu = \dfrac{\Delta N}{\Delta V} = \dfrac{N_A\Delta m /\mu}{\Delta V} = \dfrac{\rho N_A}{\mu} = 7.2\cdot10^{24}~\text{m}^{-3}, ~E \approx 9\cdot 10^4~\text{Pa}.\]

\par\smallskip\noindent\textbf{Answer:} \textbackslash{}textbf{Answer:} \[E = 3k\nu T \approx 9\cdot 10^4~\text{Pa}\]
\end{solcard}

\begin{solcard}{D1}{Write down the expression for $U_\text{терm}$. Express your answer in terms of the heat capacity at constant length $C_l$ and $T$. Consider that $C_l$ does not depend on temperature.}{0.20}
\par\smallskip\noindent\textbf{Answer:} \textbackslash{}textbf{Answer:} $$U_\text{терm} = C_lT$$
\end{solcard}

\begin{solcard}{D2}{Find an expression for $U_\text{mех}$. Express your answer in terms of $V,~\varepsilon,~T,~E(T)$ and $dE(T)/dT$.}{1.20}

As follows from the expression for entropy, during isothermal tension the entropy of the rod decreases. This means that heat is supplied on an isotherm with a higher temperature, and the Carnot cycle occurs in the forward direction. On an isotherm with a higher temperature, the rod is compressed from $\varepsilon$ to zero deformation with a constant Young's modulus $E=E(T +dT)$ and does positive work \[ A_+ = \int\limits_0^\varepsilon E\varepsilon S_0l_0\,d\varepsilon = \dfrac{E\varepsilon^2V}{2}. \] The change in internal energy on the isotherm is equal to the change in elastic strain energy, since $U_\text{терm}$ depends only on temperature: \[ \Delta U_+ =U_\text{mех}(0,T+dT) - U_\text{mех}(\varepsilon,T+dT)  = 0 - U_\text{mех}(\varepsilon,T+dT) = - U_\text{mех}.\]The heat received from the heater is \[ Q_+ = A_+ + \Delta U_+ =  \dfrac{E\varepsilon^2V}{2} - U_\text{mех}.\] By condition work on adiabats can be neglected. Then the work of the rod per cycle is equal to \[ dA =  \dfrac{E(T+dT)\varepsilon^2V}{2} -  \dfrac{E(T)\varepsilon^2V}{2} = \dfrac{ \varepsilon^2 V\,dE}{2}.\]Carnot cycle efficiency \[ \eta = \dfrac{dT}{T} = \dfrac{dA}{Q_+}. \] From here we find

\par\smallskip\noindent\textbf{Answer:} \textbackslash{}textbf{Answer:} \[ U_\text{mех} = \dfrac{\varepsilon^2 V}{2}\left(E - T\dfrac{dE}{dT}\right)\].
\end{solcard}

\begin{solcard}{D3}{Let the Young's modulus of rubber be equal to $E_0$ at temperature $T_0$. Find Young's modulus $E$ at temperature $T$ in the “ideal rubber” model. Express your answer using $E_0,~T_0$ and $T$.}{0.40}

In the “ideal rubber” model $U_\text{mех} =0$. Taking D2 into account, we have $$ E - T\dfrac{dE}{dT} = 0,$$ hence $$ \ln\dfrac{E}{E_0} = \ln\dfrac{T}{T_0}.$$


$$ E - T\dfrac{dE}{dT} = 0,$$


where


$$ \ln\dfrac{E}{E_0} = \ln\dfrac{T}{T_0}.$$

\par\smallskip\noindent\textbf{Answer:} \textbackslash{}textbf{Answer:} \[ E= E_0\dfrac{T}{T_0}\]
\end{solcard}

\begin{solcard}{E1}{Obtain an expression for $\Delta T$. Express your answer in terms of $E_0,\,S_0,\,l_0,\,C_l,\,\alpha,\,T_0$ and $\lambda$. Consider that $\Delta T = T - T_0 \ll T_0$.}{0.80}

Let's differentiate the equation of state with respect to temperature at a constant length: \textbackslash{}begin{multline} \textbackslash{}left(\textbackslash{}dfrac{\textbackslash{}partial F}{\textbackslash{}partial T}\textbackslash{}right)_l= \textbackslash{}dfrac{E_0S_0}{3T_0}\textbackslash{}left( \textbackslash{}lambda - \textbackslash{}dfrac{1+3\textbackslash{}alpha (T-T_0)}{\textbackslash{}lambda^2}\textbackslash{}right) - \textbackslash{}dfrac{E_0S_0T}{3T_0}\textbackslash{}cdot\textbackslash{}dfrac{3\textbackslash{}alpha}{\textbackslash{}lambda^2} =\textbackslash{}\textbackslash{}= \textbackslash{}dfrac{E_0S_0}{3T_0}\textbackslash{}left( \textbackslash{}lambda - \textbackslash{}dfrac{1+3\textbackslash{}alpha (T-T_0)+3\textbackslash{}alpha T}{\textbackslash{}lambda^2}\textbackslash{}right) \textbackslash{}approx \textbackslash{}dfrac{E_0S_0}{3T_0}\textbackslash{}left( \textbackslash{}lambda - \textbackslash{}dfrac{1+3\textbackslash{}alpha T_0}{\textbackslash{}lambda^2}\textbackslash{}right). \textbackslash{}end{multline}The last transition used the condition $\Delta T \ll T_0$. Substituting the resulting expression into the relation for the temperature change, we find \textbackslash{}begin{multline} \textbackslash{}Delta T = \textbackslash{}dfrac{E_0S_0}{3C_l}\textbackslash{}int\textbackslash{}limits_1^\textbackslash{}lambda\textbackslash{}left(\textbackslash{}lambda - \textbackslash{}dfrac{1+3\textbackslash{}alpha T_0}{\textbackslash{}lambda^2}\textbackslash{}right) l_0\textbackslash{},d\textbackslash{}lambda =\textbackslash{}\textbackslash{}= \textbackslash{}dfrac{E_0S_0l_0}{3C_l}\textbackslash{}left(\textbackslash{}dfrac{\textbackslash{}lambda^2 -1}{2} +(1+3\textbackslash{}alpha T_0)\textbackslash{}left(\textbackslash{}dfrac{1}{\textbackslash{}lambda} -1\textbackslash{}right)\textbackslash{}right)= \textbackslash{}dfrac{E_0S_0l_0}{3C_l}(\textbackslash{}lambda - 1)\textbackslash{}left(\textbackslash{}dfrac{\textbackslash{}lambda +1}{2} - \textbackslash{}dfrac{1+3\textbackslash{}alpha T_0}{\textbackslash{}lambda}\textbackslash{}right). \textbackslash{}end{multline}

\par\smallskip\noindent\textbf{Answer:} \textbackslash{}textbf{Answer:} \[ \Delta T = \dfrac{E_0S_0l_0}{3C_l}(\lambda - 1)\left(\dfrac{\lambda +1}{2} - \dfrac{1+3\alpha T_0}{\lambda}\right). \]
\end{solcard}

\begin{solcard}{E2}{Construct a high-quality graph $\Delta T(\lambda)$.}{0.50}

Using the resulting formula, we will construct a qualitative graph $\Delta T(\lambda)$.

\begin{figure}[H]
  \centering
  \includegraphics[width=0.72\textwidth]{T2_S_EN_assets/fig2.jpg}
\end{figure}
\end{solcard}

\begin{solcard}{E3}{At what stretching $\lambda_0$ is the sign change of $\Delta T$ observed? Give a numerical value as your answer.}{0.50}

At $\lambda = \lambda_0$ the expression in right parentheses becomes zero: \[ \dfrac{\lambda +1}{2} - \dfrac{1+3\alpha T_0}{\lambda} = 0\Rightarrow \lambda^2 +\lambda -2(1+3\alpha T_0) = 0. \] The positive root of this equation is \[ \lambda_0 = \dfrac{-1+\sqrt{1+8(1+3\alpha T_0)}}{2}\approx 1 + 2\alpha T_0 = 1.132.\]

\par\smallskip\noindent\textbf{Answer:} \textbackslash{}textbf{Answer:} $$ \lambda_0 \approx1.132$$
\end{solcard}

\begin{solcard}{E4}{What is the change in temperature $\Delta T_{1,2}$ of the sample under tension $\lambda_1= 1.1$ and $\lambda_2 = 1.4$? Use $l_0 = 1~\text{m}$ and $C_l = 150~\text{J}/\text{K}$ for calculations.}{0.40}

Using the formula from E1, we find the temperature changes:

\par\smallskip\noindent\textbf{Answer:} \textbackslash{}textbf{Answer:} \[\Delta T_1 = -1.05\cdot10^{-3}~\text{K}\]\[\Delta T_2 = 3.00\cdot10^{-2}~\text{K}\]
\end{solcard}

\end{document}